%=============================================================================
\section{Discussion and limitations}
\label{sec:limitations}
%=============================================================================

\paragraph{Assumptions and conditions of failure.}
The projections rest on three assumptions, each of which can be tested once
Spanish data exist. \emph{Directional stability}: the sign and ordering of the
effects linking gender, age, recycling behaviour, and environmental attitudes
to acceptance hold across populations. The projections fail if these effects
reverse in Spain. On ISSP proxy outcomes, effects transfer from Finland to
Spain with a slope of %%TA0%%, but the gender effect is weaker in Spain, so
the gender contrast is the least secure part of the profile ordering.
\emph{Level transfer for the reference profile}: the base rate $p_0$ is taken
unchanged from the source population. On the same proxies, Spanish and Finnish
levels differ by %%TLEVEL%% on the probit scale, so absolute levels are the
least reliable part of the output.
\emph{Construct equivalence}: source items measure the target construct on a
comparable scale. Heterogeneity statistics describe disagreement between
entries; they do not establish equivalence.

\paragraph{State of the evidence base.}
The audit of the deposited priors identifies limitations that go beyond small
samples and that bound every projection in Section~\ref{sec:scenarios}.
(i)~Each target outcome rests on one or two studies. (ii)~Pooled sample sizes
count the same respondents several times, so effective sample sizes overstate
the information available. (iii)~Some extracted entries did not measure the construct to which they were
assigned, risk items were pooled without reversal, a single percentage was
expanded into five categories without preserving it, and some statistics were
entered twice. These have been corrected through an explicit rules table
(\texttt{evidence\_rules.csv}, 15 entries), applied before pooling. They are
also the reason only one outcome is projected, and external evidence no longer
enters observed variables (Section~\ref{sec:step3}). All outputs of the
reported run are deposited (Data availability).

\paragraph{Uncertainty that is not represented.}
Intervals reflect parameter and structural uncertainty and, in the transport
specification, cross-national error measured on proxy outcomes; whether the
target outcome transports like its proxies cannot be checked. Summing pseudo-counts
narrows the prior as entries accumulate and does not represent disagreement
between studies; no between-study variance component is modelled; and
extraction, harmonisation, and transfer uncertainty are omitted. The reported
intervals are lower bounds. The power-prior discount of
Section~\ref{sec:step3} limits the weight of the transferred evidence when
Spanish data arrive but does not widen the projections themselves. The transport
error of Section~\ref{sec:transport} is a first empirical version of a
hierarchical prior over $p_0$ and the coefficients; a full version would need
the target outcome itself measured in more than one country.

\paragraph{Block-wise structure learning.}
Partitioning 90 polytomous variables into blocks makes structure learning
tractable, an order of magnitude beyond earlier applications
\citep{martins2024}. The cost is the attenuation of most cross-block
dependencies to between a fifth and two fifths of their observed strength
(Table~\ref{tab:vratios}). Queries that fully specify a profile are less
affected than queries that specify demographics alone, which are answered
through the attenuated paths and understate demographic differences. The
projections of Section~\ref{sec:scenarios} are not affected, because they are
computed from the CIS respondents' own characteristics. Three remedies have
not been tested here: a demographic super-block; overlapping partitions with
local model averaging \citep{lu2012lsbn}; and enumeration of all networks
within a Bayes-factor bound of the best \citep{liao2020credible}, which would
replace sampled graphs with a credible set. Comparing these with the present
sampler, together with Bayesian-network baselines in
Table~\ref{tab:comparison}, is required before claims about cross-block
fidelity can be generalised.

\paragraph{Associations, not effects.}
The surface describes conditional associations in a model. It does not show
that demographic characteristics cause acceptance differences, nor that any
flanking measure would change acceptance. Assessing interventions requires
experimental or quasi-experimental variation that the present design does not
contain.

\paragraph{Stated acceptance and behaviour.}
The model concerns stated acceptance. Stated preferences frequently exceed
revealed willingness to pay, in particular for unfamiliar products or products
associated with risk, so the projections should not be read as forecasts of
purchases. Linking BVS priors to scanner data or discrete-choice experiments
would allow the gap to be estimated for specific product categories.

\paragraph{Missing data and survey design.}
Mode imputation of the network variables (Section~\ref{sec:data}) treats
missing responses as typical ones. MCAR is rejected in 4 of 9 tested blocks,
so the imputed distributions, and the complete-case estimates of the transfer
model, are unbiased only under assumptions that have not been verified.
The projections of Section~\ref{sec:scenarios} do not depend on the
imputation: they are computed on CIS respondents with complete predictors.
Imputation affects the synthetic population and the fidelity results. EM
estimation of the conditional tables or multiple imputation within the
structure sampler, with results combined across imputations, would represent
this uncertainty. The treatment of survey weights is noted in Section~\ref{sec:data}.

\paragraph{Relation to other modelling approaches.}
BVS sits between generative models, which optimise fidelity to one table
\citep{xu2019}, and causal inference, which requires identification
assumptions and interventional data unavailable here. It extends
Bayesian-network synthetic data \citep{martins2024} to outcomes absent from the
observed dataset through an explicit transfer model. The baselines in
Table~\ref{tab:comparison} cannot represent such outcomes at all; the transfer
component is the part of BVS that they lack, and it should be judged by its
assumptions and, once Spanish data exist, by its predictions. Agent-based models
generate heterogeneous populations but need behavioural rules specified in
advance; acceptance surfaces could supply such rules once validated.

\paragraph{Further work.}
Beyond the corrections and extensions above, three steps follow: validation of
the target outcomes against a dedicated Spanish survey; an evaluation of the
interface in two parts, parsing accuracy and latency on a labelled set of
queries coded by hand, and a study with regulators and firms of how they read
intervals and status labels and whether the answers change their decisions;
and longitudinal data to model how acceptance moves as the
mandate takes effect.

%=============================================================================
\section{Conclusion}
\label{sec:conclusion}
%=============================================================================

Regulation~2025/40 fixes recycled content in food-contact packaging before the
consumer response to it has been measured \citep{eu2025}. We have argued that
this is a demand-side Collingridge dilemma, derived from anticipatory
governance and risk-mitigation theory what an instrument for that window must
do, and presented BVS as an AI-enabled decision-support system built to those
requirements: language-model-assisted evidence extraction, a Bayesian
inference engine that generates Spanish survey microdata and attaches an
outcome the survey does not measure through a documented transfer model, and a
natural-language interface.

The evaluation separates three levels of support. The system reproduces
variables observed in national data (mean TVD %%TVDMEAN%% over 78 variables)
but attenuates dependencies across thematic blocks. External evidence, which
did not transfer to observed Spanish variables, is confined to the outcome the
national survey does not measure. The projections for that outcome, including a 75-point difference in
projected food-packaging acceptance between contrasting profiles that ranges
from 55 to 83 points across transfer factors, are scenarios conditional on
transferred evidence and remain unvalidated in Spain.

For governance, the value of such a system lies in making assumptions
explicit, showing where adoption risk may concentrate, and directing pilots,
engagement, and measurement before commitments harden. It does not substitute
for fieldwork, for public participation, or for the accountability of those who
decide. As mandates continue to precede behavioural evidence across
circular-economy policy \citep{molenbuur2025}, instruments of this kind will be
used; the conditions under which they should and should not be used need to be
specified with the same care as their algorithms.

